% Zumdahl Chapter 16 free-response set -- 2 long (10) + 3 short (4) = 32 pts.
\documentclass{apchem-exam}

\apcunitword{Chapter}
\apcunit{16}
\apcunittitle{Solubility Equilibria}
\apcdoctitle{Free-Response Set}
\apcsubtitle{Zumdahl \S16.1--\S16.2 \textbullet\ 5 questions \textbullet\ 32 points \textbullet\ 50 minutes}

\begin{document}
\apctitleblock

\begin{apcnote}[Scope --- re-verified against the 2024 CED]
\S16.1 is fully examinable: \textbf{CED 7.11} (solubility equilibria),
\textbf{7.12} (common-ion effect) and \textbf{8.11} (pH and solubility).

\S16.2 is \textbf{partially} examinable. Predicting whether a precipitate
forms by comparing $Q$ with $K_{sp}$ is supported by the framework and is
examined in Question 2. Qualitative-analysis schemes are not.

\S16.3 (complex ion equilibria) is \textbf{off the framework entirely} ---
the phrase ``complex ion'' appears nowhere in the CED --- and no question
uses it.

One exclusion shapes Question 2: under topic 8.11, pH and solubility is
assessed \textbf{qualitatively only}. You will be asked to predict a
direction and justify it, never to compute a solubility from a pH.
\end{apcnote}

\examsection{II}{Free Response}{5 questions \textbullet\ 32 points \textbullet\ 50 minutes}{\calculatorok}

% =====================================================================
\frq{long}{10}{Zumdahl \S16.1 \textbullet\ CED 7.11, 7.12 \textbullet\ \emph{K}\textsubscript{sp}, molar solubility and the common ion}

Silver chromate, \ce{Ag2CrO4}, has $K_{sp} = 1.1\times10^{-12}$ at
\SI{25}{\celsius}.

\begin{parts}
  \item Write the dissolution equation and the $K_{sp}$ expression.
        \answerlines[2]{\ce{Ag2CrO4(s) <=> 2Ag+(aq) + CrO4^2-(aq)}

        $K_{sp} = [\ce{Ag+}]^2[\ce{CrO4^2-}]$}
  \item Calculate the molar solubility in pure water.
        \workspace[3.4cm]
        \keyonly{\begin{apcanswer}Let $s$ be the molar solubility. Then
        $[\ce{CrO4^2-}] = s$ and $[\ce{Ag+}] = \mathbf{2}s$.

        $K_{sp} = (2s)^2(s) = 4s^3 = 1.1\times10^{-12}$

        $s^3 = 2.75\times10^{-13} \;\Rightarrow\;
        s = \mathbf{6.5\times10^{-5}}$~M

        The factor of 2 is the whole difficulty here: writing
        $K_{sp} = s^3$ gives $1.0\times10^{-4}$, which is
        wrong.\end{apcanswer}}
  \item Calculate the concentration of silver ion in that saturated
        solution. \answerlines[2]{$[\ce{Ag+}] = 2s = 2(6.5\times10^{-5})
        = \mathbf{1.3\times10^{-4}}$~M}
  \item Predict, with justification, the effect on the molar solubility of
        dissolving the solid in \SI{0.10}{\Molar} \ce{AgNO3} instead of
        pure water. \workspace[3cm]
        \keyonly{\begin{apcanswer}The solubility \textbf{falls}
        substantially.

        \ce{AgNO3} supplies \ce{Ag+}, a \textbf{common ion} and a product
        of the dissolution equilibrium. Raising its concentration shifts
        the equilibrium \emph{left} by Le Ch\^{a}telier's principle, so
        less solid dissolves.

        $K_{sp}$ itself is \textbf{unchanged} --- it depends only on
        temperature. What changed is how the fixed product of
        concentrations is distributed between the two
        ions.\end{apcanswer}}
\end{parts}

\begin{rubric}
  \rpoint{2}{(a) 1 pt the equation; 1 pt the expression with \ce{Ag+}
    squared}
  \rpoint{3}{(b) 1 pt $[\ce{Ag+}] = 2s$; 1 pt the $4s^3$ form; 1 pt
    $s = \SI{6.5e-5}{\Molar}$. Using $K_{sp} = s^3$ caps this part at 1}
  \rpoint{2}{(c) 1 pt doubling $s$; 1 pt \SI{1.3e-4}{\Molar}}
  \rpoint{3}{(d) 1 pt solubility decreases; 1 pt naming the common ion and
    the leftward shift; 1 pt $K_{sp}$ unchanged. Saying $K_{sp}$ decreases
    earns \textbf{0} for that point}
\end{rubric}

% =====================================================================
\frq{long}{10}{Zumdahl \S16.1--\S16.2 \textbullet\ CED 7.11, 8.11 \textbullet\ will a precipitate form?}

Equal volumes of \SI{2.0e-3}{\Molar} \ce{AgNO3} and
\SI{2.0e-3}{\Molar} \ce{Na2CrO4} are mixed.

\begin{parts}
  \item State the concentration of each ion immediately after mixing,
        before any reaction. \answerlines[2]{Mixing equal volumes halves
        each concentration: $[\ce{Ag+}] = \SI{1.0e-3}{\Molar}$ and
        $[\ce{CrO4^2-}] = \SI{1.0e-3}{\Molar}$.}
  \item Calculate $Q$ and determine whether a precipitate of
        \ce{Ag2CrO4} forms. \workspace[3.2cm]
        \keyonly{\begin{apcanswer}$Q = [\ce{Ag+}]^2[\ce{CrO4^2-}]
        = (1.0\times10^{-3})^2(1.0\times10^{-3}) =
        \mathbf{1.0\times10^{-9}}$

        Compare with $K_{sp} = 1.1\times10^{-12}$.

        $Q > K_{sp}$ --- by about three orders of magnitude --- so the
        solution is supersaturated and a precipitate \textbf{does}
        form. Solid deposits until $Q$ falls back to
        $K_{sp}$.\end{apcanswer}}
  \item State the general rule relating $Q$ and $K_{sp}$ to what is
        observed. \answerlines[3]{$Q < K_{sp}$: unsaturated, no
        precipitate, and more solid would dissolve. $Q = K_{sp}$:
        saturated and at equilibrium. $Q > K_{sp}$: supersaturated, so
        solid precipitates until equality is restored.}
  \item Calcium carbonate is far more soluble in acidic solution than in
        neutral water, but silver chloride is essentially unaffected by
        pH. Explain the difference. \emph{Do not calculate.}
        \workspace[3.2cm]
        \keyonly{\begin{apcanswer}\ce{CaCO3} releases \ce{CO3^2-}, the
        conjugate base of the \emph{weak} acid \ce{H2CO3}. Added \ce{H+}
        reacts with carbonate and removes it from solution, so the
        dissolution equilibrium shifts \emph{right} and more solid
        dissolves.

        \ce{AgCl} releases \ce{Cl-}, which comes from the \emph{strong}
        acid \ce{HCl} and is therefore a negligible base. It does not
        react appreciably with \ce{H+}, nothing is removed from solution,
        and the equilibrium does not shift.

        The rule: a salt's solubility is pH-sensitive when its anion is
        the conjugate base of a weak acid.\end{apcanswer}}
\end{parts}

\begin{rubric}
  \rpoint{2}{(a) 1 pt recognizing dilution on mixing; 1 pt both
    \SI{1.0e-3}{\Molar}. Using the undiluted \SI{2.0e-3}{\Molar} caps the
    whole question at 6}
  \rpoint{3}{(b) 1 pt $Q$ expression with \ce{Ag+} squared; 1 pt
    $Q = \SI{1.0e-9}{}$; 1 pt the comparison and conclusion}
  \rpoint{2}{(c) 1 pt all three cases stated; 1 pt correctly linked to
    what is observed}
  \rpoint{3}{(d) 1 pt carbonate is the conjugate base of a weak acid;
    1 pt \ce{H+} removes it so the equilibrium shifts right; 1 pt chloride
    is a negligible base so no shift occurs. Note the CED assesses this
    \textbf{qualitatively only} --- do not award or expect a calculation}
\end{rubric}

% =====================================================================
\frq{short}{4}{Zumdahl \S16.1 \textbullet\ CED 7.11 \textbullet\ from solubility back to \emph{K}\textsubscript{sp}}

The molar solubility of \ce{BaSO4} in water at \SI{25}{\celsius} is
measured as \SI{1.05e-5}{\Molar}.

\begin{parts}
  \item Calculate $K_{sp}$. \answerlines[2]{\ce{BaSO4(s) <=> Ba^2+ +
        SO4^2-}, so both ions equal $s$:

        $K_{sp} = s^2 = (1.05\times10^{-5})^2 =
        \mathbf{1.1\times10^{-10}}$}
  \item Explain why this salt's expression is $s^2$ while silver
        chromate's is $4s^3$. \answerlines[3]{\ce{BaSO4} dissolves in a
        $1:1$ ratio, so each ion has concentration $s$ and neither is
        raised to a power above one. \ce{Ag2CrO4} gives \emph{two}
        \ce{Ag+} per formula unit, so its silver concentration is $2s$ and
        is squared in the expression, producing $(2s)^2(s) = 4s^3$.}
\end{parts}

\begin{rubric}
  \rpoint{2}{(a) 1 pt recognizing the $1:1$ stoichiometry; 1 pt
    \SI{1.1e-10}{}}
  \rpoint{2}{(b) 1 pt the $1:1$ versus $2:1$ stoichiometry; 1 pt the
    resulting powers}
\end{rubric}

% =====================================================================
\frq{short}{4}{Zumdahl \S16.1 \textbullet\ CED 7.11 \textbullet\ comparing solubilities}

\ce{AgCl} has $K_{sp} = 1.8\times10^{-10}$ and \ce{Ag2CrO4} has
$K_{sp} = 1.1\times10^{-12}$.

\begin{parts}
  \item A student concludes that \ce{Ag2CrO4} is less soluble because its
        $K_{sp}$ is smaller. Evaluate this reasoning.
        \workspace[3cm]
        \keyonly{\begin{apcanswer}The reasoning is \textbf{invalid}. The
        two salts have different dissolution stoichiometries, so their
        $K_{sp}$ values are not directly comparable.

        \ce{AgCl}: $s = \sqrt{1.8\times10^{-10}} =
        \SI{1.3e-5}{\Molar}$.

        \ce{Ag2CrO4}: $s = \sqrt[3]{1.1\times10^{-12}/4} =
        \SI{6.5e-5}{\Molar}$.

        The chromate is in fact about \textbf{five times more}
        soluble, despite the smaller $K_{sp}$. Comparing $K_{sp}$ directly
        is only valid for salts of the same
        stoichiometric type.\end{apcanswer}}
\end{parts}

\begin{rubric}
  \rpoint{4}{2 pts stating the reasoning is invalid because the
    stoichiometries differ; 1 pt both molar solubilities computed; 1 pt
    the correct conclusion that \ce{Ag2CrO4} is the more soluble. A
    candidate who agrees with the student earns at most 1, and only for
    correct arithmetic}
\end{rubric}

% =====================================================================
\frq{short}{4}{Zumdahl \S16.1 \textbullet\ CED 8.11 \textbullet\ pH and solubility, qualitatively}

For each solid, predict whether its solubility increases, decreases or is
essentially unchanged when the solution is acidified, and justify in one
sentence.

\begin{parts}
  \item \ce{Mg(OH)2} \answerlines[2]{\textbf{Increases.} \ce{OH-} is
        removed by reaction with \ce{H+} to form water, so the
        dissolution equilibrium shifts right.}
  \item \ce{CaF2} \answerlines[2]{\textbf{Increases.} Fluoride is the
        conjugate base of the weak acid \ce{HF}, so \ce{H+} removes it
        from solution.}
  \item \ce{AgNO3}-derived \ce{AgBr} \answerlines[2]{\textbf{Essentially
        unchanged.} Bromide comes from the strong acid \ce{HBr} and is a
        negligible base, so it does not react with \ce{H+}.}
\end{parts}

\begin{rubric}
  \rpoint{3}{1 pt each correct prediction}
  \rpoint{1}{justifications that consistently turn on whether the anion is
    the conjugate base of a weak acid. Answers of the form ``acid
    dissolves things'' earn 0 for this point}
\end{rubric}

\examtotal

\end{document}
